\documentclass{article}

\usepackage{neurips_2025}

\usepackage[utf8]{inputenc}
\usepackage[T1]{fontenc}
\usepackage{hyperref}
\usepackage{url}
\usepackage{booktabs}
\usepackage{amsfonts}
\usepackage{amsmath}
\usepackage{amssymb}
\usepackage{nicefrac}
\usepackage{microtype}
\usepackage{xcolor}
\usepackage{multirow}

\title{Variational Free Energy as a Generalized Utility Function\\ for $\rho$-POMDPs: Principled Epistemic Foraging\\ via Active Inference}

\author{%
  Anonymous Author(s)
}

\begin{document}

\maketitle

\begin{abstract}
We investigate whether Expected Free Energy (EFE), cast as the belief-dependent utility $\rho$ in $\rho$-POMDPs, produces principled epistemic foraging without per-environment parameter tuning. We compare the resulting VFE agent against five baselines---myopic reward maximization, same-horizon reward-only planning, untuned and per-environment tuned information gain, and a \texttt{pymdp} reference active inference agent---across five discrete POMDP environments. VFE's advantage is both environment- and horizon-dependent: on simple observe-then-commit problems, same-horizon reward-only planning matches VFE; on multi-observation-action environments requiring the agent to choose \emph{which} information to gather, VFE significantly outperforms same-horizon planning ($+9.1$ pp on sequential diagnosis at $H=3$, $+20.6$ pp on a structured bandit; $p < 0.05$). Per-environment tuned information gain achieves high accuracy but requires weights varying from $w=20$ to $w=100$ and systematically over-explores. VFE achieves competitive or superior performance across all environments with zero tuning, identifying structured uncertainty with sufficient planning depth as the regime where EFE-as-$\rho$ is most valuable.
\end{abstract}

\section{Introduction}

Decision-making under partial observability requires agents to balance exploiting current knowledge against gathering information to reduce uncertainty about hidden states. In standard POMDPs, information gathering has no intrinsic value; it is only instrumentally useful insofar as it leads to higher expected reward. This creates a well-known difficulty: the exploration--exploitation trade-off must be resolved either by the planning horizon or by heuristic exploration bonuses.

The $\rho$-POMDP framework \citep{araya2010} addresses this by extending POMDPs with a belief-dependent utility $\rho(b)$ that allows the agent to derive value directly from properties of its belief state. This enables explicit optimization over uncertainty reduction, information gain, or other belief-state properties alongside task reward. However, the choice of $\rho$ remains largely heuristic and environment-specific.

Separately, the Active Inference (AIF) framework \citep{friston2010, parr2019} casts perception and action as approximate Bayesian inference, selecting policies that minimize Expected Free Energy (EFE). EFE naturally decomposes into a pragmatic term (goal-seeking) and an epistemic term (information-seeking), with both arising from the same variational objective without tunable weights.

We propose substituting EFE as $\rho$ in $\rho$-POMDPs, yielding an agent whose epistemic foraging is a consequence of its objective rather than an engineered bonus. Our contributions are:
\begin{enumerate}
    \item A formal bridge between $\rho$-POMDPs and active inference, embedding EFE within the $\rho$-POMDP architecture via recursive tree search derived from sophisticated inference \citep{friston2021}.
    \item Controlled comparison against five baselines---including a same-horizon reward-only planner that isolates the effect of EFE from planning depth, and per-environment tuned information gain that gives the strongest possible alternative.
    \item Characterization of the regimes where EFE-as-$\rho$ provides genuine advantage: multi-observation-action environments with structured uncertainty and sufficient planning depth.
\end{enumerate}

\section{Related work}

\paragraph{POMDPs and belief-space planning.}
POMDPs provide the standard formalism for sequential decision-making under state uncertainty \citep{kaelbling1998}. Exact solutions over the continuous belief simplex are PSPACE-complete; point-based methods \citep{shani2013} approximate the piecewise-linear-convex value function on reachable beliefs. A persistent limitation is that information gathering has no intrinsic value: exploration is justified only instrumentally through eventual reward, making myopic solvers systematically under-explore.

\paragraph{$\rho$-POMDPs.}
\citet{araya2010} introduced $\rho$-POMDPs, augmenting the reward with a belief-dependent utility $\rho : \Delta(S) \to \mathbb{R}$, so the objective becomes $\max_\pi \mathbb{E}_\pi[\sum_t \gamma^t (R(s_t,a_t) + \rho(b_t))]$. When $\rho$ is convex, the value function retains PWLC structure, preserving compatibility with standard solvers. \citet{fehr2018} extended this to Lipschitz-continuous non-convex $\rho$, and \citet{lim2023} developed $\rho$-POMCPOW for continuous-space $\rho$-POMDPs. Common choices for $\rho$---entropy reduction, KL divergence from a target belief, information gain---each encode a different notion of epistemic value, but the choice remains heuristic and environment-specific.

\paragraph{Active inference and EFE.}
The Active Inference framework \citep{friston2010} casts perception and action as variational inference. \citet{dacosta2020} synthesized discrete-state AIF from first principles, showing that the posterior over policies takes the form $Q(\pi) \propto \exp(-\mathcal{G}(\pi))$ where $\mathcal{G}$ is the Expected Free Energy. \citet{parr2019} showed that EFE decomposes into pragmatic value (divergence from preferred observations) and epistemic value (expected information gain), with both arising from a single variational bound---no tunable exploration weight. Despite different constructions, EFE and Generalised Free Energy produce identical policy posteriors.

\paragraph{Sophisticated inference.}
Standard AIF evaluates policies myopically. \citet{friston2021} introduced \emph{sophisticated inference}, a recursive extension implementing deep tree search over belief trajectories rather than states. Sophistication---maintaining beliefs about future beliefs---enables counterfactual reasoning about the downstream epistemic consequences of actions. \citet{dacosta2020b} proved that this recursive scheme recovers Bellman-optimal policies for any finite horizon, whereas standard AIF achieves optimality only for single-step planning. Our recursive VFE agent (Equation~\ref{eq:efe_recursive}) is derived from this framework, adapted to the $\rho$-POMDP commit-action structure.

\paragraph{Exploration in RL.}
Bayes-Adaptive MDPs \citep{duff2002, guez2013} address model uncertainty rather than state uncertainty. Intrinsic motivation methods \citep{schmidhuber1991, oudeyer2007, pathak2017, bellemare2016} drive agents toward novel states via prediction error or pseudo-counts, operating on model parameters rather than belief states. VIME \citep{houthooft2016} computes information gain about dynamics via Bayesian neural networks but requires a tunable bonus weight. All these approaches engineer exploration as an auxiliary objective; our work derives it from EFE.

\paragraph{Bridging AIF and RL.}
\citet{sajid2021} demonstrated that behaviors typically engineered into RL agents---directed exploration, uncertainty-sensitive decisions, reward-free learning via preference learning---arise naturally under EFE on FrozenLake, establishing AIF as a viable alternative to RL. \citet{millidge2020} showed formal equivalence between AIF and control-as-inference under specific generative model choices. Our work takes AIF's viability as established and asks a different question: what happens when EFE is used as $\rho$ within the $\rho$-POMDP formalism? This reframing enables controlled comparison of different $\rho$ functions within the same architecture, inherits formal properties of $\rho$-POMDPs (convexity, Lipschitz continuity), and yields a more nuanced evaluation than prior AIF-vs-RL comparisons.

\section{Methodology}

\subsection{The $\rho$-POMDP framework}

A $\rho$-POMDP extends the standard POMDP with a belief-dependent utility $\rho : \Delta(S) \rightarrow \mathbb{R}$. The agent's objective becomes:
\begin{equation}
    \pi^* = \arg\max_\pi \mathbb{E}_\pi \left[ \sum_{t=0}^{H} \gamma^t \left( R(s_t, a_t) + \rho(b_t) \right) \right]
\end{equation}
When $\rho = 0$, we recover the standard POMDP. When $\rho$ encodes information gain, the agent is explicitly rewarded for reducing uncertainty.

\subsection{Expected Free Energy as $\rho$}

We define $\rho$ via the EFE from active inference. For a policy $\pi$ at future time $\tau$:
\begin{equation}
    \mathcal{G}(\pi) = \underbrace{- \mathbb{E}_{Q(o_\tau|\pi)}[\ln P(o_\tau|C)]}_{\text{Pragmatic value}} - \underbrace{\mathbb{E}_{Q(o_\tau|\pi)}\left[D_{\mathrm{KL}}\left[Q(s_\tau|o_\tau, \pi) \,\|\, Q(s_\tau|\pi)\right]\right]}_{\text{Epistemic value}}
\end{equation}
where $Q(s_\tau|\pi)$ and $Q(s_\tau|o_\tau,\pi)$ are the prior and posterior beliefs, $P(o_\tau|C)$ encodes preferred outcomes, and $D_{\mathrm{KL}}$ measures expected information gain. The pragmatic term drives goal-seeking; the epistemic term drives uncertainty reduction. Both emerge from the same variational bound with no tunable weight.

Following the sophisticated inference scheme of \citet{friston2021}, our VFE agent evaluates actions recursively:
\begin{equation}
    \mathcal{G}(\text{observe}_k) = c_k - I_k(b) + \mathbb{E}_{o}\!\left[\min_a \mathcal{G}(a \mid b'_o)\right]
    \label{eq:efe_recursive}
\end{equation}
where $c_k$ is the cost of observation action $k$, $I_k(b)$ is its expected information gain, and the expectation recurses over posterior beliefs up to planning horizon $H$. For commit actions, $\mathcal{G}(\text{commit}_i) = -\mathbb{E}_b[R_i]$.

\subsection{Agents}

We compare five agents sharing the same $\rho$-POMDP architecture, differing only in objective and planning depth:

\begin{table}[h]
\centering
\caption{Agent specifications. All use exact Bayesian belief updates.}
\label{tab:agents}
\begin{small}
\begin{tabular}{lccl}
\toprule
Agent & $\rho$ function & Horizon & Key property \\
\midrule
Myopic & $\rho = 0$ & $H=1$ & Weakest baseline \\
Planning & $\rho = 0$ & $H > 1$ & Controls for planning depth \\
Info Gain & $w \cdot \Delta H(b)$ & $H=1$ & Additive epistemic bonus \\
Info Gain (tuned) & $w^* \cdot \Delta H(b)$ & $H=1$ & Per-env.\ weight via grid search \\
VFE & EFE (Eq.~\ref{eq:efe_recursive}) & $H > 1$ & No tunable weight \\
\bottomrule
\end{tabular}
\end{small}
\end{table}

A sixth agent---a wrapper around the \texttt{pymdp} library \citep{heins2022}---validates mathematical consistency of our EFE computation against an independent implementation.

\section{Experiments}

All environments follow an observe-then-commit structure implemented as OpenAI Gymnasium environments. The agent selects among observation actions (gathering noisy information) and commit actions (terminating the episode). All experiments use 1{,}000 episodes (500 for navigation/scaling); statistical significance is assessed via independent two-sample $t$-tests. Full environment specifications are in Appendix~\ref{app:envs}.

\paragraph{Tiger} \citep{kaelbling1998}. Two states, one observation action (listen, accuracy 0.85), two commit actions. Rewards: correct $+10$, incorrect $-100$, listen cost $-1$. Tests whether agents gather sufficient evidence under extreme reward asymmetry.

\paragraph{Two-state testbed.} Two states, one observation action (accuracy 0.75), two commit actions. Rewards: correct $+1$, incorrect $-1$, observation cost $-0.1$. Isolates over-exploration under mild penalty.

\paragraph{Sequential diagnosis.} $N{=}4$ conditions, $K{=}2$ binary tests (accuracy 0.80), $N$ diagnose actions. Rewards: correct $+10$, incorrect $-50$, test cost $-1$. The agent must choose \emph{which} test to run.

\paragraph{Structured bandit.} $K{=}4$ arms, $K$ inspect actions (accuracy 0.80, cost $-0.5$), $K$ pull actions. Best arm $+10$, others $+1$. The agent must choose \emph{which} arm to inspect.

\paragraph{Navigation.} $3{\times}3$ grid, hidden goal, proximity-based observations. Goal reward $+20$, step cost $-0.5$. Movement both gathers information and changes position.

\section{Results}

\subsection{Tiger problem}

\begin{table}[h]
\centering
\caption{Tiger problem (1{,}000 episodes, $H{=}6$). Reward asymmetry $+10 / {-}100$.}
\label{tab:tiger}
\begin{small}
\begin{tabular}{lccc}
\toprule
Agent & Listens & Success & Reward \\
\midrule
Myopic & $1.00$ & 84.4\% & $-8.16$ \\
Planning ($H{=}6$) & $4.20$ & 99.3\% & $+5.03$ \\
Info Gain ($w{=}1$) & $1.00$ & 85.4\% & $-7.06$ \\
Info Gain ($w{=}20$) & $4.20$ & 99.6\% & $+5.36$ \\
\textbf{VFE} ($H{=}6$) & $\mathbf{4.29}$ & $\mathbf{99.5\%}$ & $\mathbf{+5.16}$ \\
PyMDP-AIF & $2.68$ & 97.5\% & $+4.57$ \\
\bottomrule
\end{tabular}
\end{small}
\end{table}

VFE and Planning at the same horizon ($H{=}6$) both achieve ${\sim}99\%$ success with positive reward; the differences are not statistically significant ($p = 0.29$ for observations, $p = 0.75$ for reward). On Tiger, multi-step reward planning alone discovers the value of listening because the reward structure penalizes premature commitment. Tuned Info Gain ($w{=}20$) performs equivalently. The implication: on simple observe-then-commit environments, VFE achieves parity with tuned alternatives without per-environment weight selection.

\subsection{Two-state testbed}

\begin{table}[h]
\centering
\caption{Information-seeking testbed (1{,}000 episodes, $H{=}4$). Symmetric rewards $+1 / {-}1$.}
\label{tab:infoseek}
\begin{small}
\begin{tabular}{lccc}
\toprule
Agent & Obs. & Success & Reward \\
\midrule
Myopic & $1.00$ & 73.3\% & $+0.37$ \\
Planning ($H{=}4$) & $3.22$ & 89.2\% & $\mathbf{+0.46}$ \\
Info Gain ($w{=}1$) & $3.27$ & 89.1\% & $+0.46$ \\
Info Gain ($w{=}50$) & $12.08$ & 99.9\% & $-0.21$ \\
\textbf{VFE} ($H{=}4$) & $5.50$ & $\mathbf{95.6\%}$ & $+0.36$ \\
PyMDP-AIF & $3.36$ & 91.5\% & $+0.49$ \\
\bottomrule
\end{tabular}
\end{small}
\end{table}

Under mild reward asymmetry, VFE gathers more information than is instrumentally optimal ($+0.36$ vs.\ Planning's $+0.46$, $p < 0.001$) because uncertainty reduction has intrinsic value under EFE. Tuned Info Gain ($w{=}50$) over-explores catastrophically: 99.9\% success but negative reward. Whether VFE's constitutive epistemic drive is beneficial depends on the cost-of-error regime.

\subsection{Sequential diagnosis}

\begin{table}[h]
\centering
\caption{Diagnosis ($N{=}4$, $K{=}2$ tests, $H{=}3$, 1{,}000 episodes).}
\label{tab:diagnosis}
\begin{small}
\begin{tabular}{lccc}
\toprule
Agent & Tests & Success & Reward \\
\midrule
Myopic & $2.00$ & 64.2\% & $-13.48$ \\
Planning ($H{=}3$) & $5.90$ & 87.9\% & $-3.16$ \\
Info Gain ($w{=}1$) & $2.00$ & 62.4\% & $-14.56$ \\
Info Gain ($w{=}100$) & $13.24$ & 98.7\% & $-4.02$ \\
\textbf{VFE} ($H{=}3$) & $\mathbf{9.73}$ & $\mathbf{97.0\%}$ & $\mathbf{-1.53}$ \\
\bottomrule
\end{tabular}
\end{small}
\end{table}

Here, VFE significantly outperforms Planning at the same horizon: $97.0\%$ vs.\ $87.9\%$ success ($p = 0.022$). The 9.1 pp gap demonstrates that EFE's information gain term provides genuine value in multi-observation-action environments by guiding the agent toward the most informative test. VFE also achieves the best reward ($-1.53$), outperforming tuned Info Gain ($-4.02$), which over-explores with 13.24 tests.

\subsection{Structured bandit}

\begin{table}[h]
\centering
\caption{Structured bandit ($K{=}4$ arms, $H{=}2$, 1{,}000 episodes).}
\label{tab:bandit}
\begin{small}
\begin{tabular}{lccc}
\toprule
Agent & Inspections & Success & Reward \\
\midrule
Myopic & $2.09$ & 64.6\% & $+5.77$ \\
Planning ($H{=}2$) & $3.26$ & 67.5\% & $+5.44$ \\
Info Gain ($w{=}1$) & $2.04$ & 60.5\% & $+5.43$ \\
Info Gain ($w{=}50$) & $10.90$ & 99.7\% & $+4.52$ \\
\textbf{VFE} ($H{=}2$) & $\mathbf{4.94}$ & $\mathbf{88.1\%}$ & $\mathbf{+6.46}$ \\
\bottomrule
\end{tabular}
\end{small}
\end{table}

VFE achieves the highest reward ($+6.46$, $p < 0.001$ vs.\ all others) \emph{and} highest success among reward-positive agents. Planning barely improves over Myopic ($67.5\%$ vs.\ $64.6\%$, n.s.), confirming that the bandit's symmetric reward structure provides insufficient gradient for pure reward-based exploration. VFE is the only agent that simultaneously achieves high success and high reward, demonstrating the value of intrinsic epistemic drive when the agent must decide \emph{which} information to gather.

\subsection{Summary of findings}

The results reveal a consistent pattern across environments:

\begin{itemize}
    \item \textbf{Simple observe-then-commit} (Tiger, testbed): VFE, Planning, and tuned Info Gain achieve equivalent success. VFE's advantage is tuning-free robustness.
    \item \textbf{Multi-observation-action} (Diagnosis, Bandit): VFE significantly outperforms Planning at the same horizon ($p < 0.05$). EFE's epistemic term guides \emph{which} observation action to take.
    \item \textbf{Very small spatial environments} (Navigation; see Appendix~\ref{app:navigation}): VFE over-explores. A greedy agent outperforms when exhaustive search is cheap.
    \item \textbf{Tuned Info Gain}: Achieves high accuracy but over-explores, with required weights varying from $w{=}20$ to $w{=}100$, demonstrating the fragility VFE avoids.
    \item \textbf{Horizon dependence}: At $H{=}2$, VFE and Planning perform equivalently on the scaling analysis (Appendix~\ref{app:scaling}). VFE's advantage requires $H \geq 3$---sufficient depth for the information gain term to differentiate among observation actions.
\end{itemize}

\section{Discussion}

\paragraph{When does EFE-as-$\rho$ help?}
The central finding is that VFE's advantage is not uniform but emerges under specific conditions: (1)~the environment offers multiple observation actions with differential informativeness, requiring the agent to choose \emph{which} information to gather; and (2)~the planning horizon is sufficient for the recursive EFE evaluation to propagate epistemic value through the belief tree. When either condition fails---single observation action (Tiger), or shallow horizon ($H{=}2$ in the scaling analysis)---reward-only planning achieves equivalent performance because the reward structure alone provides sufficient gradient for exploration.

\paragraph{The intrinsic exploration--exploitation balance.}
EFE resolves the exploration--exploitation trade-off dynamically: at high belief entropy, the information gain term $I(b)$ dominates, favoring observation; as beliefs sharpen, $I(b)$ decays and the pragmatic commit value takes over. This transition requires no scheduling or hyperparameter adjustment. On Tiger, both VFE and Planning achieve ${\sim}99\%$ success with ${\sim}4.2$ listens, but for different reasons: VFE is driven by information gain, Planning by the downstream reward penalty. The mechanistic distinction surfaces on the Bandit, where VFE's information gain term directs inspection toward uncertain arms while Planning lacks such guidance, producing a 20.6 pp success gap.

\paragraph{Weight fragility of additive bonuses.}
Tuned Info Gain achieves 99+\% success on multiple environments but at the cost of the lowest reward, systematically over-investing in certainty. The required weight varies from $w{=}20$ (Tiger) to $w{=}100$ (Diagnosis), with no principled method for selection. A fixed weight cannot be simultaneously appropriate across environments with different reward scales. EFE avoids this because both pragmatic and epistemic terms arise from the same variational bound, scaling naturally with the environment's information structure.

\paragraph{Limitations.}
Our environments are small discrete state-spaces ($N \leq 16$). The recursive EFE evaluation has cost $\mathcal{O}(K \cdot |\mathcal{O}|^H)$, limiting practical horizons to $H{=}2$--$3$ for environments with $K \geq 3$ observation actions. The VFE agent assumes a known generative model; extending to model learning settings (where AIF and BAMDPs converge) remains future work. On very small spatial domains (3${\times}$3 grid), VFE's epistemic drive is counterproductive, suggesting that constitutive epistemics has a minimum useful problem scale. We do not yet provide a systematic rho-function landscape comparison or visualizations of the EFE decomposition dynamics, which would more fully characterize the $\rho$-POMDP design space.

\section{Conclusion}

We have shown that Expected Free Energy serves as a coherent, tractable, and empirically effective belief-dependent utility function $\rho$ for $\rho$-POMDPs. By embedding EFE within the $\rho$-POMDP architecture, we enable controlled comparison against same-horizon reward-only planning and per-environment tuned information gain---baselines that isolate the contribution of the EFE objective from planning depth and weight sensitivity respectively. The resulting picture is nuanced: VFE provides no absolute advantage on simple environments where reward structure alone incentivizes exploration, but significantly outperforms alternatives on multi-observation-action environments where the agent must choose which information to gather. The required planning depth for this advantage ($H \geq 3$) connects directly to the sophisticated inference result that recursive EFE evaluation is necessary for Bellman optimality beyond single-step horizons. Validation against \texttt{pymdp} confirms mathematical consistency. These findings identify structured uncertainty as the regime where principled epistemic foraging via EFE is most valuable, and provide a concrete bridge between the $\rho$-POMDP and active inference literatures.

\begin{ack}
Removed for anonymous submission.
\end{ack}

{\small
\bibliographystyle{plainnat}
\begin{thebibliography}{20}

\bibitem[Araya et~al.(2010)]{araya2010}
M.~Araya, O.~Buffet, V.~Thomas, and F.~Charpillet.
\newblock A {POMDP} extension with belief-dependent rewards.
\newblock In \emph{Advances in Neural Information Processing Systems}, 2010.

\bibitem[Bellemare et~al.(2016)]{bellemare2016}
M.~Bellemare, S.~Srinivasan, G.~Ostrovski, T.~Schaul, D.~Saxton, and R.~Munos.
\newblock Unifying count-based exploration and intrinsic motivation.
\newblock In \emph{Advances in Neural Information Processing Systems}, 2016.

\bibitem[Benchetrit et~al.(2025)]{lim2023}
Y.~Benchetrit, O.~Lev-Yehudi, A.~Zhitnikov, and V.~Indelman.
\newblock $\rho$-{POMCPOW}: Online planning for continuous $\rho$-{POMDP}s.
\newblock \emph{arXiv preprint arXiv:2502.02549}, 2025.

\bibitem[Da~Costa et~al.(2020)]{dacosta2020}
L.~Da~Costa, T.~Parr, N.~Sajid, S.~Veselic, V.~Neacsu, and K.~Friston.
\newblock Active inference on discrete state-spaces: A synthesis.
\newblock \emph{Journal of Mathematical Psychology}, 99:102447, 2020.

\bibitem[Da~Costa et~al.(2023)]{dacosta2020b}
L.~Da~Costa, N.~Sajid, T.~Parr, K.~Friston, and R.~Smith.
\newblock Reward maximization through discrete active inference.
\newblock \emph{Neural Computation}, 35(5):807--852, 2023.

\bibitem[Duff and Barto(2002)]{duff2002}
M.~O. Duff and A.~G. Barto.
\newblock Optimal learning: Computational procedures for {B}ayes-adaptive {M}arkov decision processes.
\newblock PhD thesis, University of Massachusetts Amherst, 2002.

\bibitem[Fehr et~al.(2018)]{fehr2018}
R.~Fehr, O.~Buffet, V.~Thomas, and J.~Dibangoye.
\newblock $\rho$-{POMDP}s have {L}ipschitz-continuous $\epsilon$-optimal value functions.
\newblock In \emph{Advances in Neural Information Processing Systems}, 2018.

\bibitem[Friston(2010)]{friston2010}
K.~Friston.
\newblock The free-energy principle: a unified brain theory?
\newblock \emph{Nature Reviews Neuroscience}, 11(2):127--138, 2010.

\bibitem[Friston et~al.(2021)]{friston2021}
K.~Friston, L.~Da~Costa, D.~Hafner, C.~Hesp, and T.~Parr.
\newblock Sophisticated inference.
\newblock \emph{Neural Computation}, 33(3):713--763, 2021.

\bibitem[Guez et~al.(2013)]{guez2013}
A.~Guez, D.~Silver, and P.~Dayan.
\newblock Scalable and efficient {B}ayes-adaptive reinforcement learning based on {M}onte-{C}arlo tree search.
\newblock \emph{Journal of Artificial Intelligence Research}, 48:841--883, 2013.

\bibitem[Heins et~al.(2022)]{heins2022}
C.~Heins, B.~Millidge, D.~Demekas, B.~Klein, K.~Friston, I.~Couzin, and A.~Tschantz.
\newblock pymdp: A {P}ython library for active inference in discrete state spaces.
\newblock \emph{Journal of Open Source Software}, 7(73):4098, 2022.

\bibitem[Houthooft et~al.(2016)]{houthooft2016}
R.~Houthooft, X.~Chen, Y.~Duan, J.~Schulman, F.~De~Turck, and P.~Abbeel.
\newblock {VIME}: Variational information maximizing exploration.
\newblock In \emph{Advances in Neural Information Processing Systems}, 2016.

\bibitem[Kaelbling et~al.(1998)]{kaelbling1998}
L.~P. Kaelbling, M.~L. Littman, and A.~R. Cassandra.
\newblock Planning and acting in partially observable stochastic domains.
\newblock \emph{Artificial Intelligence}, 101(1--2):99--134, 1998.

\bibitem[Millidge et~al.(2020)]{millidge2020}
B.~Millidge, A.~Tschantz, and C.~L. Buckley.
\newblock On the relationship between active inference and control as inference.
\newblock In \emph{International Workshop on Active Inference}, pp.\ 3--11. Springer, 2020.

\bibitem[Oudeyer et~al.(2007)]{oudeyer2007}
P.-Y. Oudeyer, F.~Kaplan, and V.~V. Hafner.
\newblock Intrinsic motivation systems for autonomous mental development.
\newblock \emph{IEEE Transactions on Evolutionary Computation}, 11(2):265--286, 2007.

\bibitem[Parr and Friston(2019)]{parr2019}
T.~Parr and K.~J. Friston.
\newblock Generalised free energy and active inference.
\newblock \emph{Biological Cybernetics}, 113(5):495--513, 2019.

\bibitem[Pathak et~al.(2017)]{pathak2017}
D.~Pathak, P.~Agrawal, A.~A. Efros, and T.~Darrell.
\newblock Curiosity-driven exploration by self-supervised prediction.
\newblock In \emph{International Conference on Machine Learning}, 2017.

\bibitem[Sajid et~al.(2021)]{sajid2021}
N.~Sajid, P.~J. Ball, T.~Parr, and K.~J. Friston.
\newblock Active inference: Demystified and compared.
\newblock \emph{Neural Computation}, 33(3):674--712, 2021.

\bibitem[Schmidhuber(1991)]{schmidhuber1991}
J.~Schmidhuber.
\newblock A possibility for implementing curiosity and boredom in model-building neural controllers.
\newblock In \emph{Proc.\ International Conference on Simulation of Adaptive Behavior}, pp.\ 222--227, 1991.

\bibitem[Shani et~al.(2013)]{shani2013}
G.~Shani, J.~Pineau, and R.~Kaplow.
\newblock A survey of point-based {POMDP} solvers.
\newblock \emph{Autonomous Agents and Multi-Agent Systems}, 27(1):1--51, 2013.

\end{thebibliography}
}

%%%%%%%%%%%%%%%%%%%%%%%%%%%%%%%%%%%%%%%%%%%%%%%%%%%%%%%%%%%%
\appendix

\section{Environment specifications}
\label{app:envs}

\begin{table}[h]
\centering
\caption{Full environment specifications.}
\begin{small}
\begin{tabular}{lccccccc}
\toprule
Environment & $|S|$ & Obs.\ actions & Commit actions & Accuracy & Obs.\ cost & Correct & Incorrect \\
\midrule
Tiger & 2 & 1 & 2 & 0.85 & $-1.0$ & $+10$ & $-100$ \\
Testbed & 2 & 1 & 2 & 0.75 & $-0.1$ & $+1$ & $-1$ \\
Diagnosis ($N{=}4$) & 4 & 2 & 4 & 0.80 & $-1.0$ & $+10$ & $-50$ \\
Bandit ($K{=}4$) & 4 & 4 & 4 & 0.80 & $-0.5$ & $+10$ & $+1$ \\
Navigation & 9 & 4 (move) & 0 (implicit) & var. & $-0.5$ & $+20$ & --- \\
\bottomrule
\end{tabular}
\end{small}
\end{table}

All environments are implemented as OpenAI Gymnasium environments with standard \texttt{reset}/\texttt{step} API. The Diagnosis environment uses $K = \lceil \log_2 N \rceil$ binary tests, each providing noisy information about a different partition of the state space. The Navigation environment uses proximity-based warm/cold signals with accuracy depending on Manhattan distance to the hidden goal.

\section{Navigation results}
\label{app:navigation}

\begin{table}[h]
\centering
\caption{Navigation ($3{\times}3$ grid, 500 episodes).}
\begin{small}
\begin{tabular}{lccc}
\toprule
Agent & Steps & Success & Reward \\
\midrule
\textbf{NavMyopic} & $\mathbf{28.03}$ & $\mathbf{92.0\%}$ & $\mathbf{+4.38}$ \\
NavInfoGain & $41.13$ & 83.6\% & $-3.84$ \\
NavVFE ($H{=}2$) & $45.55$ & 82.8\% & $-6.22$ \\
\bottomrule
\end{tabular}
\end{small}
\end{table}

The greedy NavMyopic agent outperforms both NavInfoGain and NavVFE. On a $3{\times}3$ grid, movement under uniform belief naturally covers the space, making deliberate information-gathering detours wasteful. This delimits VFE's applicability: in very small state spaces where exhaustive search is cheap, EFE-based planning is counterproductive.

\section{Scaling analysis}
\label{app:scaling}

\begin{table}[h]
\centering
\caption{Scaling analysis: Diagnosis $N = 2$ to $16$ (500 episodes, $H{=}2$).}
\begin{small}
\begin{tabular}{clccc}
\toprule
$N$ & Agent & Tests & Success & Reward \\
\midrule
\multirow{4}{*}{2}  & Myopic    & 1.00 & 80.2\% & $-2.88$ \\
                     & Planning  & 2.99 & 93.8\% & $+3.29$ \\
                     & Info Gain & 1.00 & 79.0\% & $-3.60$ \\
                     & \textbf{VFE} & \textbf{2.89} & \textbf{95.2\%} & $\mathbf{+4.23}$ \\
\midrule
\multirow{4}{*}{4}  & Myopic & 2.00 & 67.0\% & $-11.80$ \\
                     & Planning & 5.96 & 88.8\% & $-2.68$ \\
                     & Info Gain & 2.00 & 65.0\% & $-13.00$ \\
                     & VFE & 5.95 & 87.6\% & $-3.39$ \\
\midrule
\multirow{4}{*}{8}  & Myopic & 3.00 & 50.0\% & $-23.00$ \\
                     & Planning & 8.73 & 84.6\% & $-7.97$ \\
                     & Info Gain & 3.00 & 53.4\% & $-20.96$ \\
                     & VFE & 9.03 & 83.8\% & $-8.75$ \\
\midrule
\multirow{4}{*}{16} & Myopic & 4.00 & 38.8\% & $-30.72$ \\
                     & Planning & 11.58 & \textbf{81.2\%} & $\mathbf{-12.86}$ \\
                     & Info Gain & 4.00 & 40.4\% & $-29.76$ \\
                     & VFE & 11.67 & 77.6\% & $-15.11$ \\
\bottomrule
\end{tabular}
\end{small}
\end{table}

At $H{=}2$, VFE and Planning perform similarly across all $N$, both substantially outperforming Myopic and untuned Info Gain. This contrasts with the $H{=}3$ Diagnosis result (Table~\ref{tab:diagnosis}), where VFE outperformed Planning by 9.1 pp, confirming that EFE's advantage requires sufficient recursive depth. VFE-over-Myopic advantage grows from $+15.0$ pp at $N{=}2$ to $+38.8$ pp at $N{=}16$, confirming that deeper planning is increasingly valuable as state spaces grow.

\section{PyMDP validation}
\label{app:pymdp}

We validated our EFE computation against the \texttt{pymdp} library \citep{heins2022} by wrapping pymdp's standard Agent class as a baseline (PyMDP-AIF in Tables~\ref{tab:tiger}--\ref{tab:infoseek}). Both implementations agree on observe-vs-commit decisions across all tested belief states (uniform, intermediate, and confident beliefs on the two-state environments). Information gain values computed by our recursive scheme and pymdp's \texttt{control} module agree qualitatively, with differences attributable to our recursive multi-step evaluation vs.\ pymdp's single-step policy enumeration.

\end{document}
